\documentclass{article}
\usepackage{amsfonts,amsmath,amssymb,amsthm}
\usepackage{hyperref} % do NOT set [ocgcolorlinks] here!
\usepackage{nameref}
\usepackage{xcolor}

\newcommand{\claimname}{}
\newtheorem*{claiminner}{\claimname}
\newenvironment{claim}[1]
 {\renewcommand{\claimname}{#1}\claiminner}
 {\endclaiminner}

\newtheorem{lemma}{Lemma}

\newenvironment{lemmaproof}
  {\renewcommand{\qedsymbol}{$\blacksquare$}\begin{proof}}
  {\end{proof}}

\title{$\aleph_{0}$ Weekly Problem}
\author{Ravi Dayabhai \and Conrad Warren}
\date{2025-02-18}

\begin{document}

\maketitle

\section*{Problem}

Define a sequence $\left( x_{n} \right)$ as follows:

\[
    x_{0} = 0,\ x_{n+1} = \sqrt{4 + 3x_{n}}
\]

\noindent
Show that this sequence is convergent and find its limit.

\section*{Solution}

\begin{claim}{Claim}
    The sequence \(\left( x_n \right)\) defined by \(x_0 = 0\) and \(x_{n+1} = \sqrt{4 + 3x_n}\) converges, and its limit is 4.
\end{claim}

\begin{proof}

  We begin by first establishing that the sequence \(\left( x_n \right)\) is bounded.

    \begin{lemma}[Boundedness]
    \label{bounded-lemma}
      For all \(n \geq 0\), then \(x_n \leq 4\).
    \end{lemma}

    \begin{lemmaproof}

      We proceed by induction on $n$. The base case is true because \(x_{0} = 0 \leq 4\). 
      For the inductive case, we assume \(x_n \leq 4\). Then,
    \[
    x_{n+1} = \sqrt{4 + 3x_n} \leq \sqrt{4 + 3 \cdot 4} = \sqrt{16} = 4.
    \]
    By induction, \(x_n \leq 4\) for all \(n \geq 0\).

    \end{lemmaproof}

  By the Monotone Convergence Theorem, since \(\left( x_n \right)\) is bounded above (as shown by Lemma \ref{bounded-lemma}) and monotonically increasing, it converges to some limit \(L\), which we know to be nonnegative.
  Notice also that $f(x) = \sqrt{4 + 3x}$ is a continuous function preserving monotonicity (importantly, when $x$ is nonnegative). So, the continuity of $f$ and  $L = \lim_{n \to \infty} x_n = \lim_{n \to \infty} x_{n+1}$ gives: 
      \[
        L = \lim_{n \to \infty} x_{n+1} = \lim_{n \to \infty} f(x_{n}) = f\left(\lim_{n \to \infty} x_{n} \right) = f(L) = \sqrt{4 + 3L}
      \]
Squaring both sides,
\[
L^2 = 4 + 3L \implies L^2 - 3L - 4 = 0 \implies (L - 4)(L + 1) = 0.
\]
Therefore, \(L = 4\) or \(L = -1\). Since \(x_n \geq 0\), \(L = 4\).
\end{proof}

\end{document}
